\section{Existing World}

\begin{frame}{Concept Hierarchies Are Built from Multiple Knowledge Sources}
\begin{columns}[T,onlytextwidth]
\begin{column}{0.48\textwidth}
\textbf{Expert and structured knowledge}
\begin{itemize}
    \item Ontologies define concepts and their relations explicitly.
    \item Manual construction is trusted, but expensive and quickly outdated.
    \item Linguistic resources such as dictionaries, WordNet, and FrameNet can be represented as formal contexts.
\end{itemize}
\end{column}
\begin{column}{0.48\textwidth}
\textbf{Data-driven knowledge}
\begin{itemize}
    \item Text corpora can be transformed into document-feature contexts.
    \item Topic models provide latent document-topic structure.
    \item \ac{fca} turns such contexts into interpretable concept lattices.
\end{itemize}
\end{column}
\end{columns}
\end{frame}

\begin{frame}{State of the Art}
\begin{itemize}
    \item Hierarchical text classification predicts labels inside a taxonomy.
    \item \ac{fca}-based hierarchy generation derives concepts directly from formal contexts.
    \item Constrained clustering injects prior knowledge through must-link, cannot-link, and hierarchical constraints.
    \item \ac{mlb} constraints were introduced for hierarchical agglomerative clustering to express merge order.
\end{itemize}

\vspace{0.5em}
\begin{block}{Research landscape}
Prior work has strong tools for data-driven \ac{fca} hierarchies and for constrained clustering, but hierarchical \ac{mlb} constraints have not been systematically integrated into \ac{fca}.
\end{block}
\end{frame}

\begin{frame}{Why FCA Is Attractive Here}
\begin{columns}[T,onlytextwidth]
\begin{column}{0.55\textwidth}
\begin{itemize}
    \item Objects are documents.
    \item Attributes are explicit features, topics, or generated constraints.
    \item Concepts group documents by shared attributes.
    \item The lattice exposes an interpretable hierarchy of extents and intents.
\end{itemize}
\end{column}
\begin{column}{0.40\textwidth}
\centering
\begin{tikzpicture}[
    node distance=0.95cm,
    concept/.style={circle,draw,minimum size=7mm,inner sep=0pt},
    arr/.style={-{Latex[length=2mm]},thick}
]
\node[concept] (top) {$G$};
\node[concept,below left=of top] (a) {};
\node[concept,below right=of top] (b) {};
\node[concept,below=1.0cm of a] (c) {};
\node[concept,below=1.0cm of b] (d) {};
\node[concept,below right=0.9cm and 0.45cm of c] (bot) {$\emptyset$};
\draw[arr] (top) -- (a);
\draw[arr] (top) -- (b);
\draw[arr] (a) -- (c);
\draw[arr] (b) -- (d);
\draw[arr] (c) -- (bot);
\draw[arr] (d) -- (bot);
\end{tikzpicture}
\end{column}
\end{columns}
\end{frame}
